\subsection{Text-to-Motion Generation and Control}
Text-to-motion models differ in how they represent and generate motion.
MotionDiffuse and the Human Motion Diffusion Model use continuous
 diffusion~\cite{zhang2022motiondiffusetextdrivenhumanmotion,tevet2022humanmotiondiffusionmodel}.
T2M-GPT quantizes motion and predicts tokens autoregressively, while
MotionGPT treats motion tokens as a language within a pretrained
language-model framework~\cite{zhang2023generating,jiang2023motiongpt}.
MoMask and MMM use iterative masked-token completion rather than
left-to-right prediction~\cite{guo2024momask,pinyoanuntapong2024mmmgenerativemaskedmotion}.
These approaches motivate evaluating prompt robustness across more than one
text-conditioning architecture.

Motion control and editing provide a related but distinct setting.
OmniControl imposes joint-level spatial constraints~\cite{xie2024omnicontrol};
SALAD uses skeleton-aware latent diffusion for generation and
editing~\cite{hong2025salad}. MotionFix learns from source motions, edit
instructions, and edited targets, and SimMotionEdit studies text-based editing
with motion-similarity prediction~\cite{athanasiou2024motionfix,li2025simmotionedittextbasedhumanmotion}.
COMPASS instead regularizes paired prompts for a pretrained generator. It does
not require an edited target motion for each altering pair, but this weaker
supervision also limits what can be concluded about motion-level edit
compliance.

\subsection{Prompt Robustness and Relation Supervision}
SATO identifies sensitivity to synonymous wording and small text
perturbations in text-to-motion generation~\cite{chen2024sato}. Its invariance
objective motivates our preserving branch. Intentional semantic edits require
a different response: stabilizing a changed action or direction can suppress
the requested control signal. Contrast sets and counterfactually augmented
data similarly distinguish changes that test a decision boundary from
variations that should leave a prediction
unchanged~\cite{gardner2020evaluatingmodelslocaldecision,kaushik2019learning}.
Our relation labels make this distinction explicit in prompt-pair training.
The main generation tests evaluate preserving rewrites; altering pairs inform
the objective and internal diagnostics rather than a benchmark with edited
motion ground truth.

\subsection{Attribution and Representation Regularization}
Gradient-based attribution describes the local dependence of a selected
score on the input, with interpretation that depends on the model and
attribution procedure~\cite{sundararajan2017axiomatic,adebayo2018sanity}.
COMPASS applies Grad~$\times$~Input to target-motion likelihood rather than a
text-only score. Its preserving objective applies stop-gradient Siamese
alignment to the resulting motion-conditioned
representations~\cite{chen2021exploring}.
Teacher matching constrains prediction drift, following the broader principle
of learning without forgetting~\cite{li2017learning}.
The contribution is the relation-conditioned use of these components, not a
new attribution estimator. Their forward objectives and effective gradient
paths are distinguished in our method and diagnostics.